\section{Geração dos relatórios}\label{sec:relatorios}

O relatório é elaborado após a finalização de uma aplicação e documenta as considerações do profissional a partir das informações obtidas no processo avaliativo. Cada relatório pertence a uma única aplicação (Seção~\ref{sec:modelagem}) e, por meio dela, fica contextualizado: identifica o aprendente, o instrumento e sua versão, o profissional, a data e as respostas registradas, de modo que a devolutiva não é tratada como um registro isolado.

Os campos destinados ao conteúdo do relatório são \textit{conclusao}, que registra a síntese elaborada pelo profissional, e \textit{observacoesProfissional}, que permite acrescentar informações complementares. As respostas dos instrumentos permanecem na entidade \textit{respostas}, separadas da interpretação registrada no relatório. Essa separação é deliberada: o sistema não converte os resultados de uma avaliação em diagnósticos ou conclusões psicopedagógicas, e a análise e a interpretação permanecem atribuições do profissional responsável pelo atendimento.

A página do relatório disponibiliza uma função de impressão, que abre o diálogo de impressão do navegador e permite imprimir o documento ou exportá-lo em formato PDF. O layout foi adaptado por meio de uma folha de estilos aplicada apenas no momento da impressão: a página é configurada no formato A4, com margens definidas; os elementos da interface, como o cabeçalho, o menu lateral e os botões, são ocultados; as observações e a conclusão, editadas em campos de texto na tela, são impressas como texto corrido completo; cada pergunta e sua resposta são mantidas na mesma página; e o documento é encerrado com um espaço para a assinatura do profissional responsável. As respostas do tipo sim/não são apresentadas por extenso. Essa abordagem não depende de bibliotecas de geração de PDF no servidor, utilizando a capacidade nativa dos navegadores.

O documento impresso inclui, na seção ``Dados da Avaliação'', o nome do instrumento, sua categoria e sua versão (por exemplo, ``v2''), como mostra a Figura~\ref{fig:relatorio_versao}. Como cada aplicação aponta para a versão do instrumento efetivamente utilizada, e não para a mais recente, o documento continua identificando corretamente o protocolo vigente no momento da avaliação mesmo após adaptações posteriores do instrumento.

\begin{figure}[H]
    \centering
    \caption{Bloco ``Dados da Avaliação'' no relatório impresso, com identificação da versão do instrumento}
    \label{fig:relatorio_versao}
    \includegraphics[width=0.60\textwidth]{relatorio_versao.png}
    \fonte{Elaborado pela autora (2026).}
\end{figure}
